
\documentclass{article}
\usepackage{colortbl}
\usepackage{makecell}
\usepackage{multirow}
\usepackage{supertabular}

\begin{document}

\newcounter{utterance}

\centering \large Interaction Transcript for game `wordle', experiment `high\_frequency\_words\_no\_clue\_no\_critic', episode 2 with Apertus{-}v1.5{-}8B.
\vspace{24pt}

{ \footnotesize  \setcounter{utterance}{1}
\setlength{\tabcolsep}{0pt}
\begin{supertabular}{c@{$\;$}|p{.15\linewidth}@{}p{.15\linewidth}p{.15\linewidth}p{.15\linewidth}p{.15\linewidth}p{.15\linewidth}}
    \# & \multicolumn{2}{c}{Player} && \multicolumn{2}{c}{Game Master} \\
    \hline

    \theutterance \stepcounter{utterance}  
    & & & \multicolumn{4}{p{0.6\linewidth}}{
        \cellcolor[rgb]{0.9,0.9,0.9}{
            \makecell[{{p{\linewidth}}}]{
                \texttt{\tiny{[P1$\langle$GM]}}
                \texttt{Vous êtes un magicien du langage qui aime deviner des mots en utilisant les règles données.} \\
\\ 
\texttt{Bienvenue à Wordle ! Vous avez six tentatives pour deviner le mot cible, un mot valide de cinq lettres minuscules. Veuillez utiliser les balises "explication :" et "proposition :" pour fournir une explication concise pour chaque proposition.} \\
\\ 
\texttt{Par exemple, si votre proposition est "année", votre réponse devrait être} \\
\texttt{explication : c'est un mot commun de cinq lettres, et je commence ma proposition avec ce mot.} \\
\texttt{proposition : année} \\
\\ 
\texttt{Après chaque proposition, votre réponse sera validée, et vous recevrez un retour indiquant quelles lettres sont correctes (vert), quelles lettres sont correctes mais à la mauvaise position (jaune), et quelles lettres sont incorrectes (rouge). Ce retour peut être utile pour déterminer quelles lettres inclure ou exclure dans votre prochaine proposition.} \\
\\ 
\texttt{Par exemple, le retour pour "année" pourrait être :} \\
\texttt{retour\_proposition : a<jaune> n<jaune> n<vert> é<jaune> e<rouge>} \\
\\ 
\texttt{L'explication doit contenir des détails sur la manière dont le retour\_proposition est utilisé pour arriver à une nouvelle proposition.} \\
\\ 
\texttt{Ne générez aucun autre texte.} \\
\\ 
\texttt{Commençons avec votre première proposition.} \\
            }
        }
    }
     \\ \\

    \theutterance \stepcounter{utterance}  
    & \multicolumn{4}{p{0.6\linewidth}}{
        \cellcolor[rgb]{0.561,0.737,0.561}{
            \makecell[{{p{\linewidth}}}]{
                \texttt{\tiny{[P1$\rangle$GM]}}
                \texttt{explication : j'utilise le retour\_proposition pour identifier les lettres présentes dans le mot cible et leur position, puis j'ajuste ma proposition en conséquence pour augmenter les chances de trouver le mot exact en cinq lettres.} \\
\texttt{proposition : arbre} \\
            }
        }
    }
    & & \\ \\

    \theutterance \stepcounter{utterance}  
    & & & \multicolumn{2}{p{0.3\linewidth}}{
        \cellcolor[rgb]{0.9,0.9,0.9}{
            \makecell[{{p{\linewidth}}}]{
                \texttt{\tiny{[GM$|$GM]}}
                \texttt{attempts = 1} \\
\texttt{target\_word = crise} \\
\texttt{guess = arbre} \\
\texttt{guess\_feedback = a<red> r<green> b<red> r<red> e<green>} \\
            }
        }
    }
    & & \\ \\

    \theutterance \stepcounter{utterance}  
    & & & \multicolumn{4}{p{0.6\linewidth}}{
        \cellcolor[rgb]{0.9,0.9,0.9}{
            \makecell[{{p{\linewidth}}}]{
                \texttt{\tiny{[P1$\langle$GM]}}
                \texttt{retour\_proposition : a<rouge> r<vert> b<rouge> r<rouge> e<vert>} \\
\\ 
\texttt{Fournissez votre réponse uniquement dans ce format.} \\
\texttt{explication : détails} \\
\texttt{proposition : mot} \\
            }
        }
    }
     \\ \\

    \theutterance \stepcounter{utterance}  
    & \multicolumn{4}{p{0.6\linewidth}}{
        \cellcolor[rgb]{0.561,0.737,0.561}{
            \makecell[{{p{\linewidth}}}]{
                \texttt{\tiny{[P1$\rangle$GM]}}
                \texttt{explication : le retour\_proposition indique que les lettres a et e sont présentes et correctes dans les positions respectives 1 et 5, la lettre r est correcte en position 2, et les lettres b sont incorrectes en positions 3. je conserve les lettres correctes et les positions et remplace la lettre incorrecte par une lettre qui n'est pas encore utilisée dans la proposition précédente pour essayer de deviner correctement le mot en cinq lettres.} \\
\texttt{proposition : arbre} \\
            }
        }
    }
    & & \\ \\

    \theutterance \stepcounter{utterance}  
    & & & \multicolumn{2}{p{0.3\linewidth}}{
        \cellcolor[rgb]{0.9,0.9,0.9}{
            \makecell[{{p{\linewidth}}}]{
                \texttt{\tiny{[GM$|$GM]}}
                \texttt{attempts = 2} \\
\texttt{target\_word = crise} \\
\texttt{guess = arbre} \\
\texttt{guess\_feedback = a<red> r<green> b<red> r<red> e<green>} \\
            }
        }
    }
    & & \\ \\

    \theutterance \stepcounter{utterance}  
    & & & \multicolumn{4}{p{0.6\linewidth}}{
        \cellcolor[rgb]{0.9,0.9,0.9}{
            \makecell[{{p{\linewidth}}}]{
                \texttt{\tiny{[P1$\langle$GM]}}
                \texttt{retour\_proposition : a<rouge> r<vert> b<rouge> r<rouge> e<vert>} \\
\\ 
\texttt{Fournissez votre réponse uniquement dans ce format.} \\
\texttt{explication : détails} \\
\texttt{proposition : mot} \\
            }
        }
    }
     \\ \\

    \theutterance \stepcounter{utterance}  
    & \multicolumn{4}{p{0.6\linewidth}}{
        \cellcolor[rgb]{0.561,0.737,0.561}{
            \makecell[{{p{\linewidth}}}]{
                \texttt{\tiny{[P1$\rangle$GM]}}
                \texttt{explication : le retour\_proposition confirme que les lettres a, r et e sont correcement présentes aux bonnes positions, et que les lettres b ne sont pas présentes. je conserve la structure du mot "arbre" et remplace les lettres a et b par des lettres différentes pour tester une combinaison conservant la position r et e, mais avec des lettres potentiellement différentes : mot.} \\
\texttt{proposition : mot} \\
            }
        }
    }
    & & \\ \\

    \theutterance \stepcounter{utterance}  
    & & & \multicolumn{2}{p{0.3\linewidth}}{
        \cellcolor[rgb]{0.9,0.9,0.9}{
            \makecell[{{p{\linewidth}}}]{
                \texttt{\tiny{[GM$|$GM]}}
                \texttt{Error: The length of the guessed word is not 5.} \\
            }
        }
    }
    & & \\ \\

    \theutterance \stepcounter{utterance}  
    & & & \multicolumn{2}{p{0.3\linewidth}}{
        \cellcolor[rgb]{0.9,0.9,0.9}{
            \makecell[{{p{\linewidth}}}]{
                \texttt{\tiny{[GM$|$GM]}}
                \texttt{game\_result = ABORT} \\
            }
        }
    }
    & & \\ \\

\end{supertabular}
}

\end{document}
